\section{Discussion}
\label{sec:discussion}

\paragraph{Ce que l'audit change au diagnostic.}
L'échec antérieur était attribué à une borne. L'audit montre que, pour une
projection vérifiée, la borne n'est jamais atteinte. Le véritable frein se
situe dans l'\emph{image} de la paramétrisation : la synthèse FO-LQI ne produit
ni les faibles valeurs de $k_p$ ni les valeurs de $b<1$ que requièrent les
points faisables trouvés par recherche directe. Une borne n'est donc
significative que rapportée à ce que la chaîne de synthèse peut produire ; un
taux de saturation ne dit rien sans la distinction entre saturation
structurelle et saturation pathologique.

\paragraph{Ce que la campagne établit.}
Sur 540 cellules préenregistrées, aucun correcteur n'est robustement faisable.
Ce résultat négatif est mesuré, non supposé : les graines stériles sont
rapportées, les intervalles de Wilson bornent le taux de réussite, et le test
scellé est resté intact faute de candidat à tester. Libérer $b$ ou $k_p$ rend
la faisabilité \emph{nominale} accessible, mais pas la robustesse : la
liberté de réglage ne suffit pas.

\paragraph{Un verrou physique.}
Le plancher de tension à rapport cyclique minimal est une propriété du
convertisseur, non du correcteur. Le seuil $V_\mathrm{min}\ge0{,}8$~V se
trouve au-dessus de ce plancher dans la plupart des cas d'incertitude, si
bien qu'une commande saturée pendant la descente $24\to1$~V ne peut pas le
respecter. Le cahier des charges demande donc implicitement au correcteur de
se désaturer à temps, exigence qui dépend de l'anti-emballement autant que
des gains. Toute comparaison future devrait examiner si ce seuil est
compatible avec $d_\mathrm{min}=0{,}02$ et la charge, avant de chercher des
gains.

\paragraph{Un angle mort du cahier des charges.}
Des correcteurs faisables au sens des six contraintes peuvent osciller en
permanence. Pour une revue orientée vers la mise en œuvre, c'est le constat
le plus lourd de conséquences : une contrainte de moyenne ne garantit pas
l'absence de cycle limite. Nous ne modifions pas les contraintes après
observation ; nous recommandons d'ajouter, dans une étude future et
préenregistrée, une contrainte sur l'amplitude d'oscillation en fenêtre
stationnaire.

\paragraph{Caractère fractionnaire.}
Les meilleurs points trouvés ont $\lambda$ et $\mu$ égaux ou très proches
de~1. Le correcteur qui en résulte est, en pratique, d'ordre entier. Rien dans
nos résultats n'attribue un gain à la fractionnalité ; l'étude ne comporte
pas le comparateur à réglage direct et budget égal qui permettrait d'en
juger.

\section{Limites et bilan de ce qui n'a pas été démontré}
\label{sec:limites}

\begin{itemize}
\item \textbf{Modèle.} La campagne a tourné sur un simulateur commuté écrit
  en Python. La validation croisée avec Simulink ne porte que sur six
  candidats et le plant nominal ; le snubber y est supposé ($R_s=10^5~\Omega$,
  $C_s=\infty$).
\item \textbf{Défaut corrigé après campagne.} Le feedforward de l'étage~2
  utilisait les paramètres du plant perturbé au lieu des valeurs nominales
  figées. Réévalués après correction, les 540 candidats retenus restent
  infaisables, mais 18 cellules se dégradent de plus de $0{,}05$ : le défaut
  favorisait les correcteurs, et certaines marges d'étage~2 sont optimistes.
  La recherche de l'étage~2 n'a pas été relancée.
\item \textbf{Projection.} L'identité FO-LQI $\to$ FO-PIDF n'est exacte qu'en
  ordre entier, pôle exact, sur le plant nominal non saturé ; la projection du
  pôle introduit à elle seule jusqu'à \ChiffreErreurProjPole{} d'erreur en
  haut de bande.
\item \textbf{Infaisabilité non prouvée.} Aucune cellule n'est robustement
  faisable, mais cela ne prouve pas qu'aucun correcteur robuste n'existe :
  les budgets sont finis, et l'ensemble d'incertitude est un choix déclaré.
\item \textbf{Ensemble d'incertitude.} Les paramètres incertains et leurs
  amplitudes ont été choisis par nous, le cadrage ne les précisant pas.
\item \textbf{Test scellé.} Il n'a pas été joué : aucune affirmation de
  robustesse n'est donc faite, même sous forme de proportion.
\item \textbf{Comparateurs.} Les références de l'état de l'art et le
  comparateur à réglage direct à budget égal ne sont pas implémentés ; aucune
  supériorité ne peut être attribuée à la fractionnalité.
\item \textbf{Diagnostics post hoc.} Le plancher de tension et les cycles
  limites ont été identifiés après observation ; ils motivent une étude
  future préenregistrée, mais ne doivent pas être lus comme des résultats
  confirmatoires.
\item \textbf{Métaheuristiques.} Les différences significatives mesurent la
  proximité de la faisabilité, non une faisabilité atteinte, et leur
  classement dépend du bras.
\end{itemize}
